\section{Introduction}\label{sec:intro}
Un chromosome humain peut compter plusieurs dizaines de millions de bases. Lorsque la séquence qui sert de référence est révisée, seules quelques régions peuvent avoir changé. Pour les repérer, on aligne les deux versions du chromosome en mettant en regard les bases qui se correspondent. La distance d’édition résume cette comparaison par le nombre minimal de modifications nécessaires pour passer d’une séquence à l’autre.

Le temps nécessaire pour obtenir cette distance dépend de ce que l’algorithme doit examiner. NvPD calcule en parallèle les colonnes d’une ligne à partir des scores de la ligne précédente. Lorsque les séquences se ressemblent, les méthodes à fronts peuvent atteindre le même score en explorant moins de positions. L’étude cherche à déterminer pour quelles longueurs et quelles divergences le parallélisme de NvPD compense le nombre de positions à traiter.

Sur le supercalculateur ROMEO, la comparaison commence par des séquences synthétiques. À longueur et à taux d’édition fixés, plusieurs paires permettent d’observer si le classement des méthodes change avec la séquence. Les révisions de chromosomes humains présentent un autre cas, où de longues régions identiques séparent parfois deux différences. Enfin, pour construire un arbre à partir de 64 génomes du SARS-CoV-2, il faut calculer une matrice de distances. Le processeur et le GPU peuvent aligner simultanément des paires différentes. L’étude mesure si une estimation préalable de la ressemblance réduit le temps de calcul de la matrice.

Pour deux séquences d’un million de bases, le calcul classique examine \(10^{12}\) couples de positions. La section suivante présente les deux parcours de cette grille, par lignes et par fronts.
